\section{Constructing the Admission Setting}
\label{sec:construction}
\label{sec:benchmarks}

The setting makes a useful inference compete with concepts drawn from
the passage for a fixed number of memory slots. We build each episode
around a bridge concept: a piece of information that links the passage
to a future answer. The episode supplies a passage, an ordered candidate
set, utility labels, and a future question with its answer. Only the
passage and candidates reach the admission policy.

\subsection{From a passage to a useful memory}

\textbf{Decoupled} requires two inference steps: recover the bridge
from the passage, then combine it with general knowledge to answer the
future question. GPT-4.1 authors a passage about a named person together
with the bridge, question and answer. The answer is excluded from the
passage and candidate set; the question connects to the passage through
the person's name. The author supplies vivid surface distractors and
neutral passage concepts as negative utility labels, with the intended
bridge as the positive label. Mechanical filters check candidate
strings, bridge and answer exclusion, and question overlap.
Deduplication removes repeated concepts, and a blind alias audit checks
that recovering the bridge still leaves an inference to the answer
(App.~\ref{app:benchmarks}).

\begin{figure*}[t]
\centering
\begin{sdabox}{Decoupled episode}
\textbf{Passage.} Jibril had always wanted to gaze up at the world's
highest outdoor observation deck. Early one morning, he took a swift
elevator to level 148 of the Burj Khalifa, joining a crowd of eager
tourists. Below him, the city sparkled with the first rays of the desert
sun. After snapping several panoramic photos, he browsed the sleek gift
shop for souvenirs.

\medskip
\textbf{Candidates, in supplied order.} elevator, desert, gift,
observation, dubai, tourists.\\
\textbf{Useful concept.} dubai. \textbf{Surface distractors.} desert,
tourists. \textbf{Neutral concepts.} elevator, gift, observation.

\medskip
\textbf{Future question.} What language was Jibril most likely to hear
from local staff during his visit?\\
\textbf{Answer.} arabic.
\end{sdabox}
\caption{A complete Fresh-2 episode, with \emph{dubai} as the useful concept.
Admission sees the passage and candidate strings; utility labels,
question and answer are shown for explanation.}
\label{fig:episode}
\end{figure*}

With two memory slots, the city in Figure~\ref{fig:episode} competes
with the elevator, gift shop and scenery.
\textbf{Explicit} instead states the useful content, while
\textbf{Compositional} requires a multi-hop answer and permits passage
anchors in the question. \textbf{Fresh-1} and \textbf{Fresh-2} repeat
the Decoupled construction with new passages and bridges absent from
the training batteries. Each battery is frozen before measurement.

The author uses the bridge and question to identify the intended memory;
the policy predicts its value from the passage and candidates alone.
Changing admission changes the remembered concepts while keeping the
future question and reader fixed.

\subsection{Controls named for what they change}

\paragraph{Mention and position.}
The \textbf{stated-bridge control} appends \emph{The passage also mentions
bridge.} The \textbf{position control} puts that sentence before
the passage. The \textbf{distractor-named control} substitutes a distractor
for the bridge in the appended sentence. Candidates, utility labels,
question and answer stay fixed, isolating the named concept and sentence
position.

\paragraph{Competition for memory.}
Our primary \textbf{statedness-balanced pool} starts from the 363
held-out Fresh-2 episodes in their frozen order. Counting from the first
episode, odd positions keep the passage unchanged, providing the
inferred-concept condition; even positions append the naming intervention's
sentence \emph{The passage also mentions bridge.} Here \emph{bridge}
is replaced by the useful concept, providing the stated-concept condition.
Every episode retains its original candidates, including the stated
distractors, and gains four unstated distractors at odd positions or
three at even positions. These are the first four or three additional
candidates, in order, from the pool with unstated distractors described
below. Both alternations follow position in the frozen order, with no
sampling. The resulting 363 episodes average 7.93 candidates; 13.3\%
of stated candidates and 12.0\% of unstated candidates are useful.
The probe question and answer, original utility labels, frozen reader and blind
judge remain the same; added candidates receive negative utility labels.
A statedness-only rule ranks near chance: AUC is 0.515 pooled and
0.516 within episodes.

The \textbf{naming-balanced pool} names four candidates in every episode,
using five sentence templates at varying positions. Named candidates
are useful 12.5\% of the time, against a 12.6\% base rate.
The pools control different cues: statedness in the statedness-balanced
pool and intervention naming in the naming-balanced pool. Their residual
regularities point in opposite directions: intervention naming favors
the useful concept in the statedness-balanced pool, whereas overall
statedness favors distractors in the naming-balanced pool.

Two further pools test competition for the same memory budget,
starting from the standard Fresh-2 pool's 4.43 candidates per episode.
The \textbf{pool with unstated distractors} adds four confusable concepts
from the bridge's semantic category, proposed by an author model and
filtered for absence from the passage and original pool. The
\textbf{enlarged pool} adds further concepts from other episodes,
reaching 12.41 candidates on average. Both additions preserve the
original passage, candidates, utility labels and future question while
increasing the alternatives that admission must rank.

\paragraph{Context, authorship and bridge category.}
\textbf{Shuffled context} replaces the passage with another episode's
passage while retaining its candidates and labels. The
\textbf{alternate-author benchmark} uses GPT-4o to write episodes with
an implied useful concept. The \textbf{entity-bridge benchmark} uses
GPT-5.6 to place a named-entity bridge among scene-object distractors.
Together these variants test dependence on the paired passage,
authorship and the kind of concept admitted
(App.~\ref{app:replications}).
